\section{Anatomy of the error}\label{sec:anatomy}
\textbf{The corner is outside the data.}
Under Bernoulli($p$) a unit's treated-neighbour fraction $g_i$ has mean $p$ and standard deviation $s=\sqrt{p(1-p)/d}$. The estimator in~\eqref{eq:reg} nevertheless evaluates its fit at $g=1$ and at $g=0$.
We measure the distance of the two corners from the bulk of the data in standard deviations,
\begin{equation}
Z_1=\frac{1-p}{s}=\sqrt{\frac{d(1-p)}{p}},\qquad Z_0=\frac{p}{s}=\sqrt{\frac{dp}{1-p}},\label{eq:z}
\end{equation}
and under a cluster design replace $d$ by $d_{\mathrm{eff}}=1/(\rho+(1-\rho)/d)$, where $\rho$ is the share of a unit's neighbours in its own cluster.
Both distances grow as $\sqrt d$, and a regression estimator is an extrapolation over them.

\textbf{Three terms.}
Write $\bar b$ for the coefficient of the best (design-weighted least-squares) fit of the noise-free response by the features, so that $\E\that\to\bar\Delta^{\top}\bar b$ for fixed $d$ as $n$ grows. Then
\begin{equation}
\that-\gate=\underbrace{\bar\Delta^{\top}\bar b-\gate}_{\text{projection bias}}+\underbrace{\bar\Delta^{\top}(\hat b_{\text{design}}-\bar b)}_{\text{design variance}}+\underbrace{\bar\Delta^{\top}(\hat b-\hat b_{\text{design}})}_{\text{idiosyncratic variance}},\label{eq:decomp}
\end{equation}
where $\hat b_{\text{design}}$ is the fit to the noise-free response in the realised assignment.

\emph{Idiosyncratic variance.} With independent noise of variance $\sigma^2$ the third term has variance $\sigma^2L/n$, where $L=\bar\Delta^{\top}M^{-1}\bar\Delta$ and $M$ is the second-moment matrix of the features.
For \est{REG-L} on a regular graph, $z_i$ and $g_i$ are independent and the slope block of $M^{-1}$ is the inverse covariance, so
\begin{equation}
L=\frac{1}{p(1-p)}+\frac{1}{s^2}=\frac1{p(1-p)}+(Z_0+Z_1)^2=\frac{1+d}{p(1-p)}.\label{eq:leverage}
\end{equation}
Its square root grows as $\sqrt d$. For irregular graphs and other feature sets we compute $L$ from the graph and the design.

\emph{Projection bias.} If the response is $\alpha+\beta z+\gamma H(g)$ as in~\eqref{eq:family} with $H(1)-H(0)=1$, the best linear fit on $(1,z,g)$ has slope $\gamma\,\mathrm{Cov}(H(g),g)/\mathrm{Var}(g)$ on $g$, and the bias is
\begin{equation}
\bar\Delta^{\top}\bar b-\gate=\gamma\left[\frac{\mathrm{Cov}(H(g),g)}{\mathrm{Var}(g)}-1\right]\ \xrightarrow{\ d\to\infty\ }\ \gamma\,[H'(p)-1].\label{eq:bias}
\end{equation}
For $H(g)=g^2$ and $g\sim\mathrm{Bin}(d,p)/d$ the covariance ratio is exactly $2p+(1-2p)/d$, so the bias is $2p-1+(1-2p)/d$: $-0.475$ at $p=\tfrac14$, $d=20$, whatever $n$.
Bias is the gap between the chord that joins the two corners and the tangent at $p$. It depends on $Z_0$ and $Z_1$ separately, because the two corners differ when $p\neq\tfrac12$, and it does not shrink with $n$.
We compute the left side of~\eqref{eq:bias} without fitting noise, by averaging the features' moments over designs, and compare it with the bias.

\emph{Design variance.} Unless the response is in the model, residuals within two hops are correlated, so the variance of $\that$ is not the iid sum above; the natural plug-in is a sandwich estimator with this range.

\textbf{Design-based estimators pay for the corner in variance.}
A unit is fully exposed to treatment with probability $p^{d_i+1}$ and to control with $(1-p)^{d_i+1}$, so each term of the \est{HT} estimator has variance of order $p^{-(d+1)}$ times the squared outcome: \est{HT} is unbiased when the exposure condition is right, and its variance is exponential in degree, against a leverage linear in $d$ for regression. Section~\ref{sec:experiments} compares the two costs.
